\section{Random and prime-degree binary curves}\label{sec:binary}

This section covers binary curves $y^2+xy=x^3+ax^2+b$ over $\FF_{2^n}$
with $n$ prime, attacked with random subspace bases that use no Frobenius
structure, whether the curve is random or Koblitz: the Weil-descent route and the
degree question behind the subexponential claims
(\S\ref{sec:theory}), the decomposition oracles, the random-binary
leaderboard rows, and the 131-bit ECC2-131 target. Many of the descent
studies were priced for ECC2K-130 because it was the programme's named
target, but their systems are built on prime-degree fields with random
subspaces and use no Frobenius structure (the systems at $n=19,21$ say so
explicitly), so they belong here. The Gröbner and $\FF_2$ engines they run
on are in \S\ref{sec:groebner}; every study is listed in
Appendix~\ref{sec:catalogue}.

\subsection{The degree question, measured}\label{sec:binary-descent}

The subexponential claims for prime-degree binary curves rest on the
first-fall-degree heuristic: that Weil-descended summation systems have a
solving degree close to their first fall degree (FFD), which stays small.
The programme measured both, on toy fields, in about thirty preregistered
studies from late September to 24 October.

\paragraph{First fall degree stays low; the solving degree does not.} Every
Weil-descended $\Sem3$ system for $n=3$--$7$ has first fall degree 3
(\art{research/notes/index-calculus/RESEARCH_FFD_MEASUREMENT.md}), and the
HKY last fall degree $d_F$ is 3 or 4 at every reachable size---$m=2$ to
$n=32$, $m=3$ at $\ell\le6$, $m=4$ at $\ell\le3$, $m=5,6$ to $N=54$
unknowns---while the plain refutation degree rises from 5 to at least 7
and random controls rise with it (\art{research/fall_degree_bounds_20261008/}).
At $m=2$ the Semaev $d_F$ sits one below the semi-regular degree at
$n\le18$, and the gap does not grow over $n=10$--$18$.

\paragraph{The refutation degree tracks $\ell$, within one of semi-regular.}
For chained $m=3$ systems at fixed surplus the refutation degree reads 5 at
$\ell=2$, 6 at $\ell=3,4$, and exactly 7 at $\ell=5$ on all four draws of
the $(10,5)$ cell, with FFD 3 throughout
(\art{research/dreg_ell_grid_20260925/},
\art{research/dreg_surplus_control_20260925/},
\art{research/dreg_degree7_ell5_20260930/}). At $m=4$ it reads 5, 6 and at
least 7 at $\ell=1,2,3$, flat in $n$ at fixed $\ell$
(\art{research/dreg_m4_chain_20261006/}). In 10 of 14 exact cells the
semi-regular $D_{\mathrm{reg}}$ equals the measured degree and in all 14 it
is within one. With a dense in-tree engine at $\ell=6$ ($n=19$, both
Koblitz curves) the $x_4$ formulation refutes at 11 and the RR formulation at
8 on every draw, 4 above the sharp first-fall bound of 7, rising one degree
per rung from $\ell=3$ (\art{research/ic_dense_ladder_20261004/}). A
registered semi-regular law then fitted none of seven discriminating cells,
with $x_4$ inside $[d_{\mathrm{reg}},d_{\mathrm{reg}}+1]$ on all 32 draws
(\art{research/ic_semiregular_law_20261005/}).

\paragraph{Closure algorithms refute earlier, but rise.} A MutantXL-style
closure refutes all 92 unsatisfiable draws at degree 3--4 where XL needs
5--7, and none of 345 satisfiable draws
(\art{research/dreg_closure_mutantxl_20261016/}). It stays at degree 4 to
$N=42$ at $m=4$, $\ell\le4$, but rises past 4 at $m=3$, $\ell=7$ on 8 of 8
draws (\art{research/dreg_closure_scale_20261016/}); the number of unknowns
is not the driver, and of two candidate explanations only ``$\ell=7$''
survives a run at $(m,\ell)=(4,6)$, $N=54$, which refuted at degree 4
(\art{research/dreg_closure_driver_20261016/},
\art{research/dreg_closure_m4l6_20261018/}). Guessing one or two variables
restores degree 4 at $(3,7)$ on 10 of 10 draws at $1/116$--$1/260$ of the
projected degree-5 cost (\art{research/dreg_hybrid_guess_20261024/}).

\paragraph{What the measurements price.} Applied to ECC2K-130's $m=4$ row,
the closure results contradict the degree-4 floor ($+7.62$ bits over rho)
and leave degree 5 as the optimistic case, $+10.32$ bits at $w=1$ and
$+45.71$ at $w=2$; the hybrid gives $+8.43$ bits if the number of guesses
stays flat in $\ell$ and $+58.27$ if it grows linearly, which only a run at
$\ell=8$ can separate (\S\ref{sec:ladder3}). The $m=4$ cost slope on this
engine is $0.985$--$1.060$ bits per unit of $n$ against the $0.25$ an exponent
lead would need, and a symmetrised refutation degree rises about one per
unit of $\ell$ (\art{research/ic_m4_exponent_audit_20260928/},
\art{research/ic_m4_head_engine_20260929/},
\art{research/ic_symmetry_lever_slope_20260929/}).

\paragraph{Levers that did not move the degree.} Decomposition trees cut
on the curve's own structure have the same sampled cut-out degree as random
sets in all 102 cells (\art{research/ic_tree_split_structured_v2_20261011/});
nonlinear factor-base labels refute nowhere through degree 6
(\art{research/nonlinear_factor_labels_20261002/}); Hamming-weight ideals
answer correctly but resolve only after 7--18 of $n$ coordinates are
assigned (\art{research/hamming_ideal_pdp_20260930/}); the one measured
subspace-shape effect on the degree was a trace artefact
(\art{research/notes/index-calculus/RESEARCH_FACTOR_BASE_SHAPE_SEARCH.md}),
and frozen shape cells were killed at a gap of zero
(\art{research/shape_search_g1_g6_20261015/}). The offensive first-fall-degree
thread found no lever within 9 bits of the $2^N$ boundary
(\art{research/notes/index-calculus/RESEARCH_DEGREE_REDUCTION.md}). A census
over $\FF_{2^7}$ found non-subfield quasi-subfield polynomials below the
HKPYY Lemma 4.1 bound, which answers a tightness question but moves no
cost (\art{research/notes/index-calculus/RESEARCH_QUASI_SUBFIELD.md}).

\paragraph{The same question in crypto-autoresearcher.} The companion
repository measured it independently (\S\ref{sec:companion}). msolve's
solving degree of the descended $m=3$ system is 7 at every $n$ from 11 to 23,
a lower bound from $n=17$ on, where the run walls on matrix size
(\art{ar:experiments/EXP-ALBIN-008/analysis.md}); the study reads this as
evidence on the bounded side in a narrow range that cannot distinguish
bounded from slowly growing. A degree-4 reuse closure refutes all 400
unsatisfiable $\Sem3$ descents at $n=19$ that plain degree-4 Macaulay
refutes none of (\art{ar:experiments/EXP-CERTBIN-060020/}), the same effect
as the MutantXL closure above. A crossbred solver's total exponent is $0.891$
$[0.869,0.914]$ against rho's $0.498$ over $n=10$--$24$
(\art{ar:experiments/EXP-ICI-001/analysis.md}); on its older range,
$n\le20$, the fit reproduces the repository's earlier $0.863$ as $0.864$.

\paragraph{Reading.} On toy fields, the evidence says the solving degree of
Weil-descended summation systems grows with $\ell$ and stays within one of
the semi-regular prediction, while the first and last fall degrees stay
at 3--4. That is the behaviour of a generic system, and it supports no
subexponential claim. It does not refute one: the sizes are small, the
closure algorithms bend the curve by a degree or two, and the $(4,6)$ and
$\ell=8$ cells that would decide the $m=4$ question have not run. The PKM
analogue over prime fields (\S\ref{sec:prime}) is in the same state.

\subsection{Decomposition oracles}

\paragraph{The pair-table floor.} On the non-Koblitz curve
$y^2+xy=x^3+zx+1$ the meet-in-the-middle pair table is the floor every
oracle is measured against. Over Gaudry's pairs-and-solve it is
$94$--$277\times$ at $\ell=8$, $2{,}474\times$ and $4{,}244\times$ at
$\ell=9,10$, and $928\times$ and $2{,}321\times$ at $\ell=11,12$, with no
disagreement on 128 targets per rung
(\art{research/binary_mitm_oracle_l8_20261015/},
\art{research/binary_mitm_oracle_l9_l10_20261015/},
\art{research/binary_mitm_oracle_l11_l12_20261016/}; the $\ell=12$ figure is
the validation sweep's recomputation from the committed per-target medians,
the study quotes $2{,}187\times$). These are wall times on a shared host and
unvalidated; the counted evidence---8{,}068 and 16{,}092 point additions
against 2.1 and 8.4 million quartic solves---carries the result. An earlier
pairs-and-root-find oracle was $926\times$ faster than triple enumeration at
$\ell=8$ and decided $\ell=12$ (11.5 billion triples) in 15\,s
(\art{research/notes/index-calculus/RESEARCH_SEMAEV_DECOMPOSITION.md}).

\paragraph{Algebraic oracles against the floor.} No algebraic arm beats the
pair table: the cheapest is $303\times$ the floor at $n=13$ and
$3.8\cdot10^4\times$ at $n=31$ for $m=3$, and the slopes in $n$ are steeper
(F4 $0.534$, F6-IC $0.583$, floor $0.207$; at $m=4$ a wide chained solver
$0.805$ against $0.333$; \art{research/ic_m3_algebraic_vs_mitm_20261014/}).

\paragraph{SAT, SMT and WDSat are enumerators.} Every SAT, SMT and WDSat
back end spends about one conflict per candidate tuple (CDCL $0.97$--$1.02$
per triple), so the best enumerate-and-root-find oracle beats every solver
by $10^2$--$10^4$
(\art{research/notes/index-calculus/RESEARCH_PDP_SAT_SMT_NAGAO_SURVEY.md}).
The SAT pipeline decides all 11 satisfiable $n=19$, $\ell=6$ corpus
instances in 17.8\,s; SMT runs at $0.44$--$2.63\times$ the native CDCL+XOR
solver (\art{research/pdp_smt_oracle_20261008/}); WDSat solves $n=5$
($47$ against $51{,}425$ conflicts with and without XOR-Gauss) and crashes
at $n=9$ on a static allocation limit, where CryptoMiniSat takes
0.4--2.6\,s (\art{research/wdsat_rr_20260920/}). Run unmodified on the 60
published $\Sem4$ instances of Trimoska et al., WDSat is the fastest solver,
2.37\,s on the unsatisfiable $n=19$, $\ell=6$ cell against 124.6\,s for
CryptoMiniSat with XOR clauses (\art{ar:experiments/EXP-ICPERF-66fd51/};
\S\ref{sec:companion}); that is a solver comparison on toy instances, with
no relation, linear-algebra or rho column. The SAT route is closed as an
exponent lead. Isogenous and random-$b$ neighbours give identical
Macaulay ranks and FFD 3 (\art{research/pdp_isogeny_probe_20261008/}).

\paragraph{The PDP oracle programme.} The point-decomposition speed-up
programme (\art{research/pdp_oracle_speedup_20261011/}), run over
$\FF_{2^{3\ell}}$, measured the oracle in steps before root finding replaced
it: a baseline (PS-0; $338$--$452$\,ns per pair across five rungs, A/A
spread $\le8.8\%$), a gate-count model of the group-law form (PS-1;
$7{,}410$--$14{,}608$ fused ternary-logic issues against
$70{,}695$--$190{,}275$, $9.5$--$13.0\times$), a scalar group-law kernel (PS-2;
$29.8$--$40.8$\,ns per pair, $0.067$--$0.123$ of baseline, 400 corpus root
sets agreeing), and a bit-sliced kernel ($50$--$62\times$ the baseline on two
hosts, about equal to the scalar kernel). The root-finding oracle then ran
$3{,}007$--$115{,}160\times$ the PS-0 baseline at $\ell=8$--$14$, at the generic
floor's slope (\art{benchmarks/20261012-pdp-root-finder/}). A split/hybrid
solving-degree study (Koblitz $b=1$) projected its closest relation phase at
$22.8\times$ strong rho, $391\times$ signed-Frobenius rho
(\art{research/pdp_split_dreg_20261011/}); with Frobenius closure the priced
$S$ fell to $0.35$--$0.54$ of the standard base but the best arm stayed at
$36\times$ strong rho at $n=11$ and $889\times$ at $n=19$
(\art{research/pdp_split_slots_20261012/}). A shape-replay switch cost more
on all 9 cells and stays off (\art{research/pdp_shape_replay_20261015/}).

\paragraph{Norm and circuit searches at $n=83$.} The 17--19 October
autolab cycles built a native norm predicate over GF$(2^{83})$: a cutoff-8
circuit cuts AND gates $60.5\%$ and clauses $32.9\%$, an affine support limit
of 2 cuts internal coordinates from $58{,}176$ to $45{,}427$, and 1{,}237 of
1{,}800 positive-dimensional cells were certified empty. Every hard search
returned \textsc{unknown} with zero verified witnesses, and a primary-source
check found the hybrid ANF/CNF approach already established in Bosphorus,
SMFF and PolyBoRi (\art{research/autolab_cycle_20261019_v1/},
\art{research/autolab_cycle_20261019_v3/}). The companion draft
\art{papers/norm-rootfinding/} reports the same: no measured speed-up.

\paragraph{Toy end-to-end.} On toy fields the fastest index calculus, a
cached Nagao arm, was $2.33\times$ plain rho, and $\Sem4$/$\Sem3$ with
CryptoMiniSat $14.5$--$14.7\times$ (\art{research/toy_ic_e2e_candidates_20260922/}).
Frobenius transport preserves zero and lift sets; a constant term $c$
changes zero counts only; neighbouring curves show no consistent F5
advantage; symmetry quotients help three summands ($0.26$--$0.31$ of the
ordered cost) but not two ($1.06$--$1.32$)
(\art{research/toy_binary_elimination_20260924/} and the toy cluster).

\subsection{Leaderboard rows}

Table~B of \art{docs/ic/LEADERBOARD.md} has five random-binary rows,
$\FF_{2^{15}}$ to $\FF_{2^{27}}$ ($\log_2 r=14$--$24.4$), one target, $m=2$,
folded pair table with walked targets: IC is $13.1$, $10.6$, $18$, $23$ and
$38.2\times$ matched rho. All five passed the validation gate
(commit \texttt{a8d8e3cd38}; \art{research/ic_validation_gate_20261024/}),
which the ledger paper, written a day earlier, records as no row passing.
The cofactor-projected base of LEAD-19 cuts three of these rows by
$1.29$, $1.58$ and $3.36\times$ in cold $S$, leaving $6.1$--$14.4\times$ to
matched rho (\art{research/lead19_projected_base/}); it is not yet a
leaderboard row. The integer-multiple closure of 23 October cuts $S$
$3$--$4\times$ against plain seeds at degree 26--28 and remains $6$--$12\times$
the framework's rho reference (\S\ref{sec:ladder3}).

\subsection{ECC2-131}

The Certicom ECC2-131 target ($\log_2 r=130.0$, rho $2^{64.83}$ GAE with the
negation map) has a cost model, gates and two falsifiers
(\art{docs/ic/CERTICOM131_IC_PROGRAM.md}), all on toy instances. With a free
oracle the $m=3$ floor is $2^{65.2}$--$2^{67.2}$, above rho from linear
algebra alone; $m=4$ is the first usable arity and needs a solving degree of
at most 3 (generous pricing) or 1 (plain). On random-$b$ curves at $m=3$,
$n=5$--$11$, the measured degree is $D=4,5,6$ at $\ell=1,2,3$, equal to the
semi-regular reference with zero spread over 60 draws, and equal to the
Koblitz $b=1$ value on 52 of 52 paired draws; extrapolated to $n=131$ it is
between about 6--7 (if it plateaus) and 39 (semi-regular). One CPU core walks
ECC2-131 rho at 8.50 million steps per second with a batched pclmul walk, and
11.46 million bit-sliced (\art{research/certicom131_rho_ic_bench_20261022/},
\art{research/certicom131_rho_parity_20261023/}); a CUDA port compiles but
has not run on a GPU. The bench's IC ladders on smaller random-binary
curves read $11$--$24\times$ matched rho; it is not an admitted row.

\paragraph{sect113r1.} The 113-bit SECG curve is weak by size alone
($2^{56.33}$ generic work); its two degree-5 isogeny routes transfer the
problem without an established speed-up, and the GHS and Semaev routes
reached no conclusion (\art{research/sect113r1-weak-curve-20261006/}).
